\clearpage
\sectionkicker{Source-backed technical notes}
\subsection*{Inference \& Serving — モデルの外側も更新された: Technical Notes}
この欄は記事本文で圧縮した一次資料上の情報を、比較・再検証しやすい形へ展開したものである。新しい外部情報は追加せず、Selection済みEvidenceのchronology、normalized claim、limitations、source URLのみを再配置する。
\medskip
\subsection*{Theme at a glance}
\addcontentsline{toc}{subsection}{Theme at a glance}
\begin{center}
\footnotesize
\begin{tabularx}{\linewidth}{@{}p{0.20\linewidth}p{0.10\linewidth}p{0.14\linewidth}X@{}}
\toprule
Artifact & Role & Type & Objective chronology \\
\midrule
vLLM July releases & PRIMARY & FRAMEWORK & 2026-07-11 (FRAMEWORK\_RELEASE); 2026-07-27 (FRAMEWORK\_RELEASE) \\
SGLang July releases & PRIMARY & FRAMEWORK & 2026-07-10 (FRAMEWORK\_RELEASE); 2026-07-25 (FRAMEWORK\_RELEASE) \\
FlashInfer July MoE/Blackwell serving releases & PRIMARY & FRAMEWORK & 2026-07-17 (FRAMEWORK\_RELEASE); 2026-07-31 (FRAMEWORK\_RELEASE) \\
July diffusion/MoE inference research & SUPPORTING & OTHER & 2026-07-13 (PAPER\_PUBLICATION); 2026-07-14 (PAPER\_PUBLICATION) \\
\bottomrule
\end{tabularx}
\end{center}
\normalsize

\begin{technicalnote}{vLLM July releases}{PRIMARY}
\begin{tabularx}{\linewidth}{@{}>{\bfseries}p{0.22\linewidth}X@{}}
Organization & vLLM \\
Artifact type & FRAMEWORK \\
Chronology & 2026-07-11 (FRAMEWORK\_RELEASE); 2026-07-27 (FRAMEWORK\_RELEASE) \\
\end{tabularx}
\smallskip
{\bfseries 一次資料から整理したtechnical points}
\begin{itemize}[leftmargin=1.5em,itemsep=0.35em]
\item \textbf{Project claim}: vLLM v0.25.0 made Model Runner V2 the default execution path for dense models and removed the legacy PagedAttention implementation.
\item \textbf{Project claim}: vLLM v0.26.0 added broad new model-family support and expanded attention, speculative-decoding, KV-offload, and DeepSeek-V4-related optimizations.
\end{itemize}
{\bfseries 読む際の境界}
\begin{itemize}[leftmargin=1.5em,itemsep=0.35em]
\item \textbf{分析上の留意点}: Performance percentages in release notes are project measurements for specific kernels/models; integration support does not by itself prove model quality or availability outside vLLM.
\end{itemize}
{\bfseries Primary source}
\begingroup\sloppy
\begin{itemize}[leftmargin=1.5em,itemsep=0.25em]
\item {\scriptsize\url{https://github.com/vllm-project/vllm/releases/tag/v0.25.0}}
\item {\scriptsize\url{https://github.com/vllm-project/vllm/releases/tag/v0.26.0}}
\end{itemize}
\endgroup
{\scriptsize\color{SurveyMuted}Source-bound record: \texttt{evidence:SP-2026-M07:series-vllm-july-18398a2366}.}
\end{technicalnote}
\medskip

\begin{technicalnote}{SGLang July releases}{PRIMARY}
\begin{tabularx}{\linewidth}{@{}>{\bfseries}p{0.22\linewidth}X@{}}
Organization & SGLang \\
Artifact type & FRAMEWORK \\
Chronology & 2026-07-10 (FRAMEWORK\_RELEASE); 2026-07-25 (FRAMEWORK\_RELEASE) \\
\end{tabularx}
\smallskip
{\bfseries 一次資料から整理したtechnical points}
\begin{itemize}[leftmargin=1.5em,itemsep=0.35em]
\item \textbf{Project claim}: SGLang v0.5.15 made its Spec V2 path default and included production tuning for newer model/hardware combinations.
\item \textbf{Project claim}: SGLang v0.5.16 introduced DSpark, a confidence-driven speculative-decoding algorithm that varies verification windows using draft confidence.
\end{itemize}
{\bfseries 読む際の境界}
\begin{itemize}[leftmargin=1.5em,itemsep=0.35em]
\item \textbf{分析上の留意点}: Reported throughput/TPS figures are project benchmarks under specific batch size, tensor-parallel, model, and GPU configurations.
\end{itemize}
{\bfseries Primary source}
\begingroup\sloppy
\begin{itemize}[leftmargin=1.5em,itemsep=0.25em]
\item {\scriptsize\url{https://github.com/sgl-project/sglang/releases/tag/v0.5.15}}
\item {\scriptsize\url{https://github.com/sgl-project/sglang/releases/tag/v0.5.16}}
\end{itemize}
\endgroup
{\scriptsize\color{SurveyMuted}Source-bound record: \texttt{evidence:SP-2026-M07:series-sglang-july-1ce78e1eff}.}
\end{technicalnote}
\medskip

\begin{technicalnote}{FlashInfer July MoE/Blackwell serving releases}{PRIMARY}
\begin{tabularx}{\linewidth}{@{}>{\bfseries}p{0.22\linewidth}X@{}}
Organization & FlashInfer \\
Artifact type & FRAMEWORK \\
Chronology & 2026-07-17 (FRAMEWORK\_RELEASE); 2026-07-31 (FRAMEWORK\_RELEASE) \\
\end{tabularx}
\smallskip
{\bfseries 一次資料から整理したtechnical points}
\begin{itemize}[leftmargin=1.5em,itemsep=0.35em]
\item \textbf{Project claim}: FlashInfer v0.6.15 shipped expert-parallel MoE support in the default install and expanded Blackwell-family serving coverage.
\item \textbf{Project claim}: FlashInfer v0.6.16 followed with MegaMoE kernels, further Blackwell model coverage, an expanded unified MoE API, and distributed-serving changes including confidential-computing support.
\end{itemize}
{\bfseries 読む際の境界}
\begin{itemize}[leftmargin=1.5em,itemsep=0.35em]
\item \textbf{分析上の留意点}: Hardware/model support is highly GPU-generation-specific, and any throughput numbers in release notes are project benchmarks rather than general serving guarantees.
\end{itemize}
{\bfseries Primary source}
\begingroup\sloppy
\begin{itemize}[leftmargin=1.5em,itemsep=0.25em]
\item {\scriptsize\url{https://github.com/flashinfer-ai/flashinfer/releases/tag/v0.6.15}}
\item {\scriptsize\url{https://github.com/flashinfer-ai/flashinfer/releases/tag/v0.6.16}}
\end{itemize}
\endgroup
{\scriptsize\color{SurveyMuted}Source-bound record: \texttt{evidence:SP-2026-M07:series-flashinfer-july-41fd93a505}.}
\end{technicalnote}
\medskip

\begin{technicalnote}{July diffusion/MoE inference research}{SUPPORTING}
\begin{tabularx}{\linewidth}{@{}>{\bfseries}p{0.22\linewidth}X@{}}
Organization & - \\
Artifact type & OTHER \\
Chronology & 2026-07-13 (PAPER\_PUBLICATION); 2026-07-14 (PAPER\_PUBLICATION) \\
\end{tabularx}
\smallskip
{\bfseries 一次資料から整理したtechnical points}
\begin{itemize}[leftmargin=1.5em,itemsep=0.35em]
\item \textbf{Author claim}: FlashDiff proposes regional execution and scheduling to reduce repeated full-region computation and communication overhead in diffusion-model serving.
\item \textbf{Author claim}: EcoSpec proposes cost-aware speculative decoding for sparse MoE models by considering expert activation/scattering in addition to token acceptance likelihood.
\end{itemize}
{\bfseries 読む際の境界}
\begin{itemize}[leftmargin=1.5em,itemsep=0.35em]
\item \textbf{分析上の留意点}: Both are author-evaluated systems papers; reported gains depend on their selected models, hardware, workloads, and baselines.
\end{itemize}
{\bfseries Primary source}
\begingroup\sloppy
\begin{itemize}[leftmargin=1.5em,itemsep=0.25em]
\item {\scriptsize\url{http://arxiv.org/abs/2607.12121v2}}
\item {\scriptsize\url{http://arxiv.org/abs/2607.12696v1}}
\end{itemize}
\endgroup
{\scriptsize\color{SurveyMuted}Source-bound record: \texttt{evidence:SP-2026-M07:series-inference-systems-july-e02ad54ec0}.}
\end{technicalnote}
\medskip
